\documentclass[12pt]{article}
\usepackage[utf8]{inputenc}
\usepackage[T1]{fontenc}
\usepackage{hyperref}
\title{Mpinda Language Resources for K-12 Curriculum Development}
\author{Compiled by Claude Sonnet 5.5 (claude-sonnet-5-5) for Andr\'as Kornai}
\date{}
\begin{document}
\maketitle

\section{Introduction}
Mpinda (ISO 639-3 \texttt{pnd}, Glottocode mpin1234) is a Bantu (Atlantic-Congo) language of Angola, described by Wikipedia as closely related to Kimbundu, with about 18,000 native speakers (2018). Joshua Project places the Mpinda mainly in southeastern Angola (Cuando Cubango, Bi\'e, Moxico, Cunene, Hu\'ila) and gives an alternate name, Mbwela, with population figures of 38,000--39,000. The figures conflict and are approximate. Nothing in the search linked the language to the Republic of the Congo. The search found no educational materials in Mpinda.

\section{Primary Educational Resources}
No primers, textbooks or school programmes in Mpinda were found.

\section{Linguistic and Reference Materials}
\begin{itemize}
\item Wikipedia, ``Mpinda language''. \url{https://en.wikipedia.org/wiki/Mpinda_language}
\end{itemize}

\section{Religious and Community Resources}
\begin{itemize}
\item Joshua Project, Mpinda in Angola: people-group profile; Bible translation listed as started, with minimal Christian resources. \url{https://joshuaproject.net/people_groups/21580/AO}
\item Joshua Project, language page for \texttt{pnd}. \url{https://joshuaproject.net/languages/pnd}
\end{itemize}

\section{Digital and Audio Resources}
No audio, video, corpus or software resources were found.

\section{Curriculum Development Context}
Mpinda is not known to be taught in schools, and no written resources were located. Angolan schooling is mainly in Portuguese. Any curriculum would require new documentation work with speaker communities.

\end{document}